%%%%%%%%%%%%%%%%%%%%%%%%%%%%%%%%%%%%%%%%%%%%%%%%%%%%%%%%%%%%%%%%%%%%%%%%
%% RELATÓRIO DESCRITIVO — gera relatorio-descritivo.pdf                %%
%%                                                                     %%
%% Uma das quatro peças do pedido de patente (art. 3º, II da            %%
%% Portaria/INPI/DIRPA nº 14/2024). Cada peça é um documento SEPARADO,  %%
%% com numeração de páginas independente (art. 16, I), e é anexada como %%
%% arquivo próprio no Peticionamento Eletrônico.                        %%
%%                                                                     %%
%% Este arquivo é só orquestração: a ordem das seções segue os incisos  %%
%% do art. 27. O texto vai nos arquivos de pedido/relatorio-descritivo/.%%
%%%%%%%%%%%%%%%%%%%%%%%%%%%%%%%%%%%%%%%%%%%%%%%%%%%%%%%%%%%%%%%%%%%%%%%%

\documentclass[a4paper,12pt]{article}
\usepackage{lib/inpitex}
%%%%%%%%%%%%%%%%%%%%%%%%%%%%%%%%%%%%%%%%%%%%%%%%%%%%%%%%%%%%%%%%%%%%%%%%
%% DADOS DO PEDIDO — o único arquivo de metadados do template.         %%
%%                                                                     %%
%% Lido pelas quatro raízes (relatorio-descritivo.tex,                 %%
%% reivindicacoes.tex, desenhos.tex e resumo.tex), de modo que o        %%
%% título seja necessariamente o mesmo no relatório descritivo e no     %%
%% resumo — como exige o art. 24, IV da Portaria/INPI/DIRPA nº 14/2024. %%
%%%%%%%%%%%%%%%%%%%%%%%%%%%%%%%%%%%%%%%%%%%%%%%%%%%%%%%%%%%%%%%%%%%%%%%%

%%%%%%%%%%%%%%%%%%%%%%%%%%%%%%%%%%%%%%%%%%%%%%%%%%%%%%%%%%%%%%%%%%%%%%%%
%% 1. TÍTULO DO PEDIDO
%%
%% Art. 24 — o título deve:
%%   I   ser conciso, claro e específico, identificando o objeto do pedido,
%%       não podendo exceder 500 caracteres;
%%   II  não conter denominações de fantasia;
%%   III não conter fórmulas químicas ou matemáticas;
%%   IV  ser o mesmo no relatório descritivo e no resumo.
%% Art. 25 — deve representar adequadamente as diferentes categorias de
%%       reivindicação (se você reivindica produto E processo, o título
%%       precisa refletir os dois).
%%%%%%%%%%%%%%%%%%%%%%%%%%%%%%%%%%%%%%%%%%%%%%%%%%%%%%%%%%%%%%%%%%%%%%%%

% Os \providecommand permitem que o showcase (./gerar-exemplos.sh) injete
% outros valores por linha de comando, sem duplicar nenhum arquivo. Na
% compilação normal valem os valores definidos aqui — é aqui que você edita.
\providecommand{\TituloDoPedido}{DISPOSITIVO DE ACIONAMENTO PARA VÁLVULA DE
IRRIGAÇÃO E PROCESSO DE CONTROLE DE VAZÃO}
\titulodopedido{\TituloDoPedido}

%%%%%%%%%%%%%%%%%%%%%%%%%%%%%%%%%%%%%%%%%%%%%%%%%%%%%%%%%%%%%%%%%%%%%%%%
%% 2. NATUREZA DA PROTEÇÃO
%%
%%   invencao             — patente de invenção. Reivindicações conforme os
%%                          arts. 29 a 31; desenhos opcionais.
%%   modelo-de-utilidade  — patente de modelo de utilidade. UMA ÚNICA
%%                          reivindicação independente (art. 33), dependentes
%%                          só nas hipóteses do art. 35, e desenhos
%%                          OBRIGATÓRIOS (art. 22).
%%
%% Se você não sabe qual é o seu caso: invenção é uma solução técnica nova
%% para um problema técnico (art. 8º da LPI); modelo de utilidade é um objeto
%% de uso prático, ou parte dele, com nova forma ou disposição que resulte em
%% melhoria funcional no seu uso ou fabricação (art. 9º da LPI).
%%%%%%%%%%%%%%%%%%%%%%%%%%%%%%%%%%%%%%%%%%%%%%%%%%%%%%%%%%%%%%%%%%%%%%%%

\providecommand{\NaturezaDoPedido}{invencao}
\natureza{\NaturezaDoPedido}

%%%%%%%%%%%%%%%%%%%%%%%%%%%%%%%%%%%%%%%%%%%%%%%%%%%%%%%%%%%%%%%%%%%%%%%%
%% 3. MODALIDADE DO DEPÓSITO
%%
%%   originario         — depósito comum. Nada é impresso após o título.
%%   adicao             — certificado de adição de invenção. Imprime, após o
%%                        título, "Certificado de adição de invenção do ___"
%%                        (art. 43, II). Só existe para natureza 'invencao'
%%                        (arts. 42 e 43) e exige \pedidovinculado.
%%   divisao            — pedido dividido. Imprime, após o título,
%%                        "Dividido do ___" (art. 51, II). Exige
%%                        \pedidovinculado.
%%   fase-nacional-pct  — entrada na fase nacional de pedido internacional
%%                        (Portaria INPI nº 39/2021). A forma dos documentos é
%%                        a mesma desta norma (art. 65); os dados do pedido
%%                        internacional vão no requerimento eletrônico, não
%%                        nos documentos.
%%%%%%%%%%%%%%%%%%%%%%%%%%%%%%%%%%%%%%%%%%%%%%%%%%%%%%%%%%%%%%%%%%%%%%%%

\providecommand{\ModalidadeDoPedido}{originario}
\modalidade{\ModalidadeDoPedido}

% Número do pedido ou da patente a que este documento se vincula. Obrigatório
% nas modalidades 'adicao' e 'divisao'; ignorado nas demais.
% Exemplo: \pedidovinculado{BR 10 2024 000000 0}
\providecommand{\PedidoVinculado}{}
\pedidovinculado{\PedidoVinculado}

%%%%%%%%%%%%%%%%%%%%%%%%%%%%%%%%%%%%%%%%%%%%%%%%%%%%%%%%%%%%%%%%%%%%%%%%
%% 4. FORMATO DO RÓTULO DOS PARÁGRAFOS
%%
%% Art. 26, II — cada parágrafo do relatório descritivo recebe numeração
%% sequencial em algarismos arábicos à esquerda do texto. A norma exemplifica
%% com três dígitos ([003], [015]); os formulários oficiais do e-Patentes usam
%% um dígito ([1], [2]). As duas formas atendem à norma.
%%
%%   tres-digitos  — [001], [002], ... (padrão)
%%   arabico       — [1], [2], ...
%%%%%%%%%%%%%%%%%%%%%%%%%%%%%%%%%%%%%%%%%%%%%%%%%%%%%%%%%%%%%%%%%%%%%%%%

\formatorotuloparagrafo{tres-digitos}

%%%%%%%%%%%%%%%%%%%%%%%%%%%%%%%%%%%%%%%%%%%%%%%%%%%%%%%%%%%%%%%%%%%%%%%%
%% 5. CÓPIA DE COMPARAÇÃO (arts. 51, III e 57, II)
%%
%% Ao apresentar modificação depois do depósito, ou ao depositar pedido
%% dividido, a norma pede dois documentos tirados da MESMA matéria: os
%% documentos modificados "sem qualquer tipo de rasura ou sinalização"
%% (art. 57, I) e uma cópia de comparação com tachado indicando remoção e
%% sublinhado indicando inclusão ou substituição (art. 57, II).
%%
%% Marque as alterações no texto com \removido{...}, \incluido{...} e
%% \substituido{...}{...} (e, para blocos inteiros, \paragraforemovido e
%% \reivindicacaoremovida). O mesmo arquivo gera as duas saídas.
%%
%% NÃO mexa no valor abaixo. Para gerar a cópia de comparação use
%%
%%     make comparacao        (ou ./gerar-copia-de-comparacao.sh)
%%
%% que produz os PDFs *-comparacao.pdf sem tocar neste arquivo. Deixar 'sim'
%% aqui faria os PDFs do pedido saírem marcados, contrariando o art. 57, I —
%% e por isso o `make verificar` recusa compilar assim.
%%%%%%%%%%%%%%%%%%%%%%%%%%%%%%%%%%%%%%%%%%%%%%%%%%%%%%%%%%%%%%%%%%%%%%%%

\providecommand{\CopiaDeComparacao}{nao}
\copiadecomparacao{\CopiaDeComparacao}

%%%%%%%%%%%%%%%%%%%%%%%%%%%%%%%%%%%%%%%%%%%%%%%%%%%%%%%%%%%%%%%%%%%%%%%%
%% 6. DADOS DO DEPÓSITO — NÃO SÃO IMPRESSOS EM NENHUM DOCUMENTO
%%
%% Guarde-os aqui para ter tudo em um lugar só, mas saiba que eles NÃO vão
%% para os PDFs: são preenchidos no formulário eletrônico de requerimento do
%% Peticionamento Eletrônico do INPI (art. 3º, I).
%%
%% O art. 21 proíbe, nos documentos do pedido, "timbres, logotipos, letreiros,
%% assinaturas ou rubricas, sinais ou indicações de qualquer natureza
%% estranhos à matéria do pedido". Nome de depositante, logotipo institucional,
%% endereço e assinatura, portanto, NÃO podem aparecer no relatório
%% descritivo, nas reivindicações, nos desenhos ou no resumo.
%%
%%   Depositante(s) ......... Empresa Brasileira de Pesquisa Agropecuária
%%   CNPJ ................... 00.348.003/0001-10
%%   Inventor(es) ........... (nomes completos, na ordem do requerimento)
%%   Procurador ............. (se houver; art. 59)
%%   Prioridade unionista ... (país, número e data do primeiro depósito; art. 12)
%%   Prioridade interna ..... (número e data do pedido anterior; art. 15)
%%   Período de graça ....... (forma, local e data da divulgação; art. 11)
%%   Código de serviço/GRU .. (natureza do pedido, paga até a apresentação)
%%   Figura que melhor representa o pedido ... Figura 1 (art. 37, § único)
%%   Listagem de sequências ... (se houver: código de controle devolvido pelo
%%                               Peticionamento ao anexar o XML no padrão ST.26;
%%                               Portaria INPI/PR nº 48/2022, arts. 4º e 6º)
%%%%%%%%%%%%%%%%%%%%%%%%%%%%%%%%%%%%%%%%%%%%%%%%%%%%%%%%%%%%%%%%%%%%%%%%


\begin{document}

% Art. 26, I — o título vem apenas na página inicial, centralizado e separado
% do texto do relatório descritivo em si. Nas modalidades 'adicao' e 'divisao',
% a menção exigida pelos arts. 43, II e 51, II sai logo abaixo do título.
\abrirrelatoriodescritivo

% Art. 27, I — precisar o setor técnico.
%%%%%%%%%%%%%%%%%%%%%%%%%%%%%%%%%%%%%%%%%%%%%%%%%%%%%%%%%%%%%%%%%%%%%%%%
%% Art. 27, I — "precisar o setor técnico a que se aplica a invenção ou %%
%% o modelo de utilidade".                                              %%
%%                                                                     %%
%% Diga a qual setor técnico o seu objeto pertence: composições de      %%
%% tintura capilar, máquinas para semeadura, comunicações de rede sem   %%
%% fio, artigos domésticos... Se o objeto puder ser aplicado em mais de %%
%% um campo técnico, cite todos.                                        %%
%%                                                                     %%
%% Abra cada parágrafo com \pnum — ele produz a numeração sequencial    %%
%% [001], [002] exigida pelo art. 26, II.                               %%
%%%%%%%%%%%%%%%%%%%%%%%%%%%%%%%%%%%%%%%%%%%%%%%%%%%%%%%%%%%%%%%%%%%%%%%%

\secaoinpi{Campo \danatureza}

\pnum A presente \nomedanatureza\ pertence ao setor técnico dos equipamentos de
irrigação agrícola, mais especificamente ao dos dispositivos de acionamento de
válvulas empregadas no controle de vazão em sistemas de irrigação localizada.

\pnum A \nomedanatureza\ aplica-se também ao setor de automação de sistemas
hidráulicos de baixa pressão, no qual o acionamento de válvulas exige resposta
rápida e consumo reduzido de energia.


% Art. 27, II a IV — estado da técnica, problema técnico, solução, vantagens,
% novidade e efeito técnico (invenção) ou melhoria funcional (modelo).
%%%%%%%%%%%%%%%%%%%%%%%%%%%%%%%%%%%%%%%%%%%%%%%%%%%%%%%%%%%%%%%%%%%%%%%%
%% Art. 27, II a IV — a seção mais decisiva para o exame:              %%
%%   II  descrever o estado da técnica útil à compreensão, à busca e ao %%
%%       exame do pedido, citando sempre que possível os documentos que %%
%%       o reflitam, e destacando os problemas técnicos existentes;      %%
%%   III descrever, de forma clara, concisa e precisa, a solução        %%
%%       proposta e as vantagens em relação ao estado da técnica;        %%
%%   IV  ressaltar nitidamente a novidade e evidenciar o efeito técnico %%
%%       alcançado (invenção) ou a melhoria funcional (modelo).          %%
%%                                                                     %%
%% Como citar o estado da técnica: em prosa, dentro dos parágrafos      %%
%% numerados, identificando o documento pelo número de publicação (ex.: %%
%% BR 10 2015 000000 0, US 2018/0123456 A1, EP 3 000 000 B1). NÃO há    %%
%% lista de referências no pedido de patente — nada de ABNT aqui. As    %%
%% diretrizes de exame (Resolução INPI/PR nº 124/2013, itens 2.40 a     %%
%% 2.46) tratam de como os documentos citados são considerados.         %%
%%                                                                     %%
%% Ponto de atenção do art. 32 da LPI, regulado pela Resolução          %%
%% INPI/PR nº 93/2013: informação essencial ao exame e à                %%
%% patenteabilidade NÃO pode ser inserida depois do requerimento de     %%
%% exame. O que não estiver aqui até lá pode custar o indeferimento,    %%
%% mesmo que o objeto seja novo e inventivo.                            %%
%%%%%%%%%%%%%%%%%%%%%%%%%%%%%%%%%%%%%%%%%%%%%%%%%%%%%%%%%%%%%%%%%%%%%%%%

\secaoinpi{Fundamentos \danatureza}

\pnum São conhecidos do estado da técnica dispositivos de acionamento de
válvulas de irrigação que empregam solenoides de acionamento direto, nos quais
a força magnética atua sobre a haste de vedação ao longo de todo o curso de
abertura. O documento BR 10 2015 000000 0 descreve um conjunto desse tipo, no
qual o solenoide permanece energizado durante todo o intervalo de irrigação.

\pnum Também são conhecidos dispositivos de acionamento hidráulico assistido,
como o descrito no documento US 2018/0123456 A1, em que a pressão da própria
linha auxilia o deslocamento do elemento de vedação, reduzindo a força exigida
do atuador.

\pnum Os dispositivos de acionamento direto do estado da técnica apresentam
consumo elevado de energia, uma vez que exigem corrente contínua de manutenção
durante todo o período de abertura da válvula, o que limita seu emprego em
sistemas alimentados por painéis fotovoltaicos de pequeno porte. Os
dispositivos de acionamento hidráulico assistido, por sua vez, dependem de uma
pressão mínima de linha para operar, e apresentam falha de fechamento quando a
pressão cai abaixo desse limite.

\pnum O problema técnico não resolvido pelo estado da técnica consiste,
portanto, em acionar a válvula com baixo consumo de energia sem tornar o
fechamento dependente da pressão da linha.

\ifinvencao
\pnum A presente invenção resolve esse problema por meio de um dispositivo de
acionamento biestável, no qual o atuador eletromecânico atua apenas durante a
transição entre os estados de abertura e de fechamento, e a mola de retorno
mantém o estado alcançado sem consumo de energia. O efeito técnico alcançado é
a redução do consumo de energia por ciclo de acionamento, com fechamento
garantido por meio mecânico, independentemente da pressão da linha.

\pnum A novidade da invenção reside na associação entre o assento de vedação de
dupla face e a mola de retorno pré-carregada, que juntos estabelecem dois
estados estáveis de operação, dispensando a corrente de manutenção exigida
pelos dispositivos de acionamento direto conhecidos.
\else
\pnum O presente modelo de utilidade resolve esse problema por meio de nova
disposição construtiva do conjunto de acionamento, na qual a haste de
acionamento recebe um ressalto de encosto que atua sobre a mola de retorno
pré-carregada, de modo que o conjunto passa a apresentar dois estados estáveis
de operação. A melhoria funcional alcançada é a redução do consumo de energia
por ciclo de acionamento, com fechamento garantido por meio mecânico,
independentemente da pressão da linha.

\pnum A nova forma reside no ressalto de encosto da haste de acionamento e em
seu posicionamento em relação ao assento de vedação de dupla face, disposição
que não se encontra nos objetos conhecidos do estado da técnica.
\fi


% Arts. 26, III e 27, V — listagem de TODAS as figuras do documento de
% desenhos. Gerada a partir de pedido/desenhos/figuras.tex, o mesmo arquivo que
% desenhos.tex usa para posicionar as imagens: numeração e ordem ficam
% consistentes por construção. Comente a seção inteira se o seu pedido não
% tiver desenhos (proibido em modelo de utilidade — art. 22).
\secaoinpi{Breve descrição dos desenhos}
\imprimirdescricaodosdesenhos

% Art. 27, VI e VII — descrição precisa, clara e suficiente, de maneira que um
% técnico no assunto possa realizá-la (art. 24 da LPI).
%%%%%%%%%%%%%%%%%%%%%%%%%%%%%%%%%%%%%%%%%%%%%%%%%%%%%%%%%%%%%%%%%%%%%%%%
%% Art. 27, VI — descrever de forma precisa, clara e suficiente, de     %%
%% maneira que um técnico no assunto possa realizá-la, fazendo remissão %%
%% aos sinais de referência constantes dos desenhos.                    %%
%%                                                                     %%
%% Esta é a maior seção do relatório descritivo. Vale o art. 24 da LPI: %%
%% tudo o que for necessário para que alguém reproduza o objeto precisa %%
%% estar AQUI. O examinador avalia a suficiência descritiva por este    %%
%% texto (Resolução INPI/PR nº 124/2013, itens 2.13 a 2.16).            %%
%%                                                                     %%
%% Sinais de referência: sempre que citar um elemento do desenho, cite  %%
%% o sinal entre parênteses — "a haste de acionamento (4) desloca o     %%
%% assento (9)". Os mesmos sinais devem identificar o mesmo elemento em %%
%% todas as figuras (art. 39, IV) e ser uniformes em todo o pedido      %%
%% (art. 23, IV).                                                       %%
%%                                                                     %%
%% Unidade do objeto:                                                   %%
%%   invenção — um único conceito inventivo (art. 22 da LPI). Motor novo %%
%%     e sistema de freios novo resolvem problemas técnicos diferentes: %%
%%     são dois pedidos.                                                %%
%%   modelo de utilidade — uma única unidade técnico-funcional e corporal%%
%%     (art. 23 da LPI). Mesa nova e cadeira nova são dois corpos com   %%
%%     funções distintas: são dois pedidos.                             %%
%%                                                                     %%
%% Atenção (art. 27, VII): em modelo de utilidade NÃO cabem trechos do  %%
%% tipo "concretização preferida" ou "a título exemplificativo" — o     %%
%% texto deve definir o objeto requerido e suas variantes, não um       %%
%% princípio que admita formas diversas.                                %%
%%%%%%%%%%%%%%%%%%%%%%%%%%%%%%%%%%%%%%%%%%%%%%%%%%%%%%%%%%%%%%%%%%%%%%%%

\secaoinpi{Descrição \danatureza}

\pnum O dispositivo de acionamento (1) compreende um corpo tubular (2), provido
de bocal de entrada e de bocal de saída, no interior do qual se aloja uma
membrana elástica (3) solidária a uma haste de acionamento (4).

\pnum A haste de acionamento (4) apresenta, em sua porção intermediária, um
ressalto de encosto que atua sobre uma mola de retorno (5) pré-carregada,
apoiada contra a face interna do corpo tubular (2). Na extremidade oposta à
membrana elástica (3), a haste de acionamento (4) recebe um assento de vedação
de dupla face (9), que veda alternadamente o bocal de saída e um orifício de
alívio do corpo tubular (2).

\pnum Um atuador eletromecânico (6) acopla-se à haste de acionamento (4) e
promove o deslocamento desta entre duas posições extremas. A pré-carga da mola
de retorno (5) é dimensionada de modo que, ultrapassado o ponto de inversão
definido pelo ressalto de encosto, a mola de retorno (5) conduza a haste de
acionamento (4) até a posição extrema seguinte e a mantenha nessa posição sem
alimentação do atuador eletromecânico (6).

\pnum Um sensor de pressão (7), montado no bocal de saída do corpo tubular (2),
fornece o sinal de pressão de linha a uma unidade de controle (8), que comanda
o atuador eletromecânico (6) por meio de pulsos de corrente de duração
determinada.

\pnum No estado de abertura, o assento de vedação de dupla face (9) mantém
vedado o orifício de alívio e liberado o bocal de saída, e a água escoa da
entrada para a saída do corpo tubular (2). No estado de fechamento, o assento
de vedação de dupla face (9) mantém vedado o bocal de saída e liberado o
orifício de alívio, o que despressuriza a câmara sob a membrana elástica (3) e
consolida o fechamento.

\pnum A transição entre os dois estados ocorre exclusivamente durante o pulso
de corrente aplicado ao atuador eletromecânico (6), cuja duração é inferior a
duzentos milissegundos, de modo que o consumo de energia por ciclo de
acionamento se reduz à energia do pulso.

\pnum O processo de controle de vazão executado pela unidade de controle (8)
compreende as etapas de: leitura do sinal de pressão fornecido pelo sensor de
pressão (7); comparação do valor lido com uma faixa de operação previamente
definida; e aplicação de um pulso de corrente ao atuador eletromecânico (6)
quando o valor lido permanecer fora da faixa de operação por intervalo superior
a um tempo de confirmação.

\pnum A etapa de comparação emprega uma faixa de operação delimitada por um
limite inferior e um limite superior de pressão, e o tempo de confirmação
impede que oscilações transitórias da linha provoquem acionamentos
desnecessários.


% Art. 27, VI — exemplos e/ou quadros comparativos. É aqui que entram as
% tabelas (art. 20); figuras, nunca (art. 19).
%%%%%%%%%%%%%%%%%%%%%%%%%%%%%%%%%%%%%%%%%%%%%%%%%%%%%%%%%%%%%%%%%%%%%%%%
%% Art. 27, VI — "se necessário, utilizar exemplos e/ou quadros        %%
%% comparativos, relacionando-os com o estado da técnica".              %%
%%                                                                     %%
%% Resultados de teste comparativo costumam ser decisivos para          %%
%% demonstrar o efeito técnico, e por isso pesam no exame. Se você tem  %%
%% esses dados, este é o lugar deles.                                   %%
%%                                                                     %%
%% TABELAS aqui, FIGURAS nunca:                                         %%
%%   art. 20 — tabelas, fórmulas ou estruturas químicas e expressões    %%
%%     matemáticas podem ser inseridas no texto, em preto e            %%
%%     identificadas de forma sequencial. Use o ambiente 'tabelainpi'   %%
%%     (numera e rotula sozinho) e o ambiente 'equation' para           %%
%%     expressões matemáticas;                                          %%
%%   art. 19 — representações gráficas (figuras, fotografias,          %%
%%     fluxogramas, gráficos) NÃO são aceitas no relatório descritivo.  %%
%%     Elas vão exclusivamente no documento de desenhos.                %%
%%                                                                     %%
%% Toda tabela precisa ser comentada no texto — uma tabela solta, sem   %%
%% parágrafo que a explique, é matéria que o examinador não sabe como   %%
%% ler.                                                                 %%
%%%%%%%%%%%%%%%%%%%%%%%%%%%%%%%%%%%%%%%%%%%%%%%%%%%%%%%%%%%%%%%%%%%%%%%%

\secaoinpi{Exemplos de concretizações \danatureza}

\ifinvencao
\pnum Em uma concretização preferida, o corpo tubular (2) é obtido por injeção
de poliamida reforçada com fibra de vidro, a membrana elástica (3) é de
elastômero de etileno-propileno-dieno e a mola de retorno (5) é de aço inoxidável
austenítico, com pré-carga de 12 N.
\else
\pnum Em uma primeira variação construtiva, o corpo tubular (2) é obtido por
injeção de poliamida reforçada com fibra de vidro, a membrana elástica (3) é de
elastômero de etileno-propileno-dieno e a mola de retorno (5) é de aço inoxidável
austenítico, com pré-carga de 12 N.
\fi

\pnum O atuador eletromecânico (6) é constituído por um solenoide de bobina
única, alimentado em 12 V, e o sensor de pressão (7) é do tipo piezorresistivo,
com faixa de medição de 0 kPa a 600 kPa.

\pnum A energia consumida por ciclo de acionamento é obtida pela integral da
potência aplicada ao atuador eletromecânico (6) ao longo da duração do pulso,
conforme a expressão:

\begin{equation}
    W = \int_{0}^{t_{p}} v(t)\, i(t)\, \mathrm{d}t
\end{equation}

\pnum Foram realizados ensaios comparativos entre uma concretização do
dispositivo de acionamento (1) e dois conjuntos representativos do estado da
técnica, submetidos ao mesmo regime de 40 ciclos diários de irrigação, com
pressão de linha de 200 kPa. Os resultados estão reunidos na Tabela 1.

\begin{tabelainpi}{Energia consumida por ciclo de acionamento e pressão mínima
de fechamento}
    \begin{tabular}{lrr}
        \toprule
        Conjunto ensaiado & Energia por ciclo (J) & Pressão mínima (kPa) \\
        \midrule
        Acionamento direto (BR 10 2015 000000 0)      & 216{,}0 & 0   \\
        Acionamento hidráulico (US 2018/0123456 A1)   & 4{,}8   & 120 \\
        Dispositivo de acionamento (1)                & 1{,}9   & 0   \\
        \bottomrule
    \end{tabular}
\end{tabelainpi}

\pnum Os valores reunidos na Tabela 1 mostram que o dispositivo de acionamento
(1) consumiu 1,9 J por ciclo, contra 216,0 J do conjunto de acionamento direto
e 4,8 J do conjunto de acionamento hidráulico assistido. O conjunto de
acionamento hidráulico assistido deixou de completar o fechamento com pressão
de linha inferior a 120 kPa, enquanto o dispositivo de acionamento (1) completou
o fechamento em todos os ensaios, inclusive com a linha despressurizada.


% Art. 27, VIII — aplicação industrial, quando não for evidente a partir da
% descrição. Esvazie o arquivo se for evidente no seu caso.
%%%%%%%%%%%%%%%%%%%%%%%%%%%%%%%%%%%%%%%%%%%%%%%%%%%%%%%%%%%%%%%%%%%%%%%%
%% Art. 27, VIII — "indicar, explicitamente, a aplicação industrial     %%
%% quando essa não for evidente a partir da descrição da invenção ou do %%
%% modelo de utilidade".                                                %%
%%                                                                     %%
%% Só é obrigatório quando a aplicação industrial NÃO for evidente. Se  %%
%% no seu caso ela for evidente, apague o conteúdo deste arquivo e      %%
%% comente a linha correspondente em relatorio-descritivo.tex.          %%
%%                                                                     %%
%% Art. 27, IX — a ordem dos incisos I a VIII pode ser alterada quando  %%
%% outra ordem permitir melhor compreensão e apresentação mais concisa. %%
%%%%%%%%%%%%%%%%%%%%%%%%%%%%%%%%%%%%%%%%%%%%%%%%%%%%%%%%%%%%%%%%%%%%%%%%

\secaoinpi{Aplicação industrial}

\pnum O dispositivo de acionamento (1) é fabricável por processos convencionais
de injeção de termoplásticos e de montagem de conjuntos hidráulicos, e destina-se
à produção em escala industrial por fabricantes de componentes para irrigação.


\end{document}
